% CORRECTED VERSION
% Changes:
% - Fixed misuse of unbiasedness in Step 2 by handling the correlation between η_t and g_t.
% - Replaced the invalid inequality in Step 5 with a correct condition α ≤ ε/(2L) and a proper bound.
% - Provided a rigorous proof of the logarithmic bound (18) in Lemma 3 and used it in Step 13.
% - Adjusted constants throughout and updated the theorem statement accordingly.

\documentclass[11pt]{article}
\usepackage{amsmath,amssymb,amsthm}
\usepackage{algorithm}
\usepackage[noend]{algpseudocode}
\usepackage{enumitem}
\usepackage{hyperref}
\usepackage{geometry}
\geometry{margin=1in}
\newtheorem{theorem}{Theorem}
\newtheorem{lemma}{Lemma}
\newcommand{\E}{\mathbb{E}}
\newcommand{\R}{\mathbb{R}}
\newcommand{\norm}[1]{\left\|#1\right\|}
\newcommand{\ip}[2]{\left\langle #1,#2 \right\rangle}
\begin{document}

\section*{RMSProp (Scalar Version)}

We analyse the simplest scalar version of RMSProp, i.e. the parameter $w\in\R$ is a single real number.  The multi‑dimensional case follows by applying the scalar proof component‑wise.

\section*{Mathematical Formulation}
\begin{align}
    &\text{Initialize } w_0\in\R,\quad v_0=0, \quad \beta\in[0,1),\;\alpha>0,\;\varepsilon>0. \nonumber\\[4pt]
    &\text{For } t=0,1,2,\dots,T-1 \text{ do: } \nonumber\\
    &\qquad g_t \;=\; \nabla f(w_t;\xi_t) \quad\text{(stochastic gradient)}\label{eq:gt}\\
    &\qquad v_{t+1} \;=\; \beta\,v_t + (1-\beta)\,g_t^{2} \label{eq:vt}\\
    &\qquad \eta_t \;=\; \frac{\alpha}{\sqrt{v_{t+1}}+\varepsilon} \label{eq:stepsize}\\
    &\qquad w_{t+1} \;=\; w_t - \eta_t\, g_t . \label{eq:update}
\end{align}

\section*{Pseudocode}
\begin{algorithm}[H]
\caption{RMSProp (scalar)}\label{alg:rmsprop}
\begin{algorithmic}[1]
\State \textbf{Input:} $w_0$, base step size $\alpha>0$, decay $\beta\in[0,1)$, $\varepsilon>0$, max‑it $T$.
\State $v_0\gets 0$
\For{$t=0$ \textbf{to} $T-1$}
    \State Sample a data point $\xi_t$ and compute stochastic gradient $g_t=\nabla f(w_t;\xi_t)$
    \State $v_{t+1}\gets \beta\,v_t + (1-\beta)g_t^{2}$
    \State $\eta_t \gets \alpha/(\sqrt{v_{t+1}}+\varepsilon)$
    \State $w_{t+1}\gets w_t - \eta_t g_t$
\EndFor
\State \textbf{Return} $w_T$ (or the averaged iterate $\bar w_T=\frac1T\sum_{t=0}^{T-1} w_t$)
\end{algorithmic}
\end{algorithm}

\section*{Assumptions}
\begin{enumerate}[label={\bf A\arabic*:}, leftmargin=*]
\item \textbf{Smoothness.} $f:\R\to\R$ is $L$‑smooth, i.e.
\[
|f(w')-f(w)-\nabla f(w)(w'-w)|\le \frac{L}{2}(w'-w)^2,\qquad\forall w,w'\in\R. \tag{A1}
\]

\item \textbf{Unbiased stochastic gradients and bounded variance.}
\[
\E[g_t\,|\,w_t]=\nabla f(w_t),\qquad 
\E\big[(g_t-\nabla f(w_t))^{2}\,\big|\,w_t\big]\le\sigma^{2}. \tag{A2}
\]

\item \textbf{Bounded second moment of true gradients.}
There exists $G>0$ such that
\[
\|\nabla f(w)\|\le G,\qquad\forall w\in\R. \tag{A3}
\]

\item \textbf{Hyper‑parameter constraints.}
\[
0\le\beta<1,\qquad \alpha>0,\qquad \varepsilon>0,\qquad 
\alpha\le \frac{\varepsilon}{2L}. \tag{A4}
\]
% ADDED: the extra condition $\alpha\le\varepsilon/(2L)$ guarantees the inequality used in Step 5.

\item \textbf{Existence of a global minimum.}
Denote $f^{*}= \inf_{w}f(w)$ (finite). \tag{A5}
\end{enumerate}

\section*{Main Convergence Theorem}
\begin{theorem}[Non‑convex RMSProp convergence]\label{thm:main}
Under assumptions \textbf{A1}–\textbf{A5} and with step‑size base $\alpha$ satisfying \textbf{A4}, the iterates generated by \eqref{eq:gt}–\eqref{eq:update} satisfy
\[
\frac{1}{T}\sum_{t=0}^{T-1}\E\big[\|\nabla f(w_t)\|^{2}\big]\;\le\;
\frac{2\big(f(w_0)-f^{*}\big)}{\alpha\,T}
\;+\;
\frac{2L\alpha\sigma^{2}}{1-\beta}\,
\frac{\log\!\big(1+\tfrac{(1-\beta)G^{2}T}{\varepsilon^{2}}\big)}{T}. \tag{1}
\]
Consequently,
\[
\min_{0\le t<T}\E\big[\|\nabla f(w_t)\|^{2}\big]\;=\;O\!\bigg(\frac{\log T}{\sqrt{T}}\bigg).
\]
\end{theorem}

\section*{Preliminaries (Definitions and Lemmas)}
\begin{lemma}[Recursive bound on $v_t$]\label{lem:vtbound}
For all $t\ge0$,
\[
v_{t+1}\;=\;(1-\beta)\sum_{k=0}^{t}\beta^{\,t-k} g_k^{2}
\;\le\;(1-\beta)\max_{0\le k\le t} g_k^{2}.
\]
Moreover, using \textbf{A3},
\[
v_{t+1}\;\le\;(1-\beta)G^{2}\,\frac{1-\beta^{t+1}}{1-\beta}
\;\le\;G^{2}. \tag{2}
\]
\end{lemma}
\begin{proof}
Unrolling the recursion \eqref{eq:vt} yields the first equality.  The remaining inequalities follow from $g_k^{2}\le G^{2}$ (A3). \qedhere
\end{proof}

\begin{lemma}[Lower bound on stepsize]\label{lem:eta}
For any $t$,
\[
\eta_t \;\ge\; \frac{\alpha}{\sqrt{1-\beta}\,|g_t|+\varepsilon}
\;=\; \frac{\alpha}{\sqrt{1-\beta}\,|g_t|+\varepsilon}. \tag{3}
\]
\end{lemma}
\begin{proof}
From \eqref{eq:vt} we have $v_{t+1}\ge (1-\beta)g_t^{2}$, hence
$\sqrt{v_{t+1}}+\varepsilon\le \sqrt{(1-\beta)}|g_t|+\varepsilon$, which gives the claimed lower bound after inversion. \qedhere
\end{proof}

\begin{lemma}[Logarithmic bound on the sum of adaptive stepsizes]\label{lem:logsum}
Define $s_t:=v_{t+1}+\varepsilon^{2}$.  Then for any $T\ge1$,
\[
\sum_{t=0}^{T-1}\eta_t^{2}
= \frac{\alpha^{2}}{1-\beta}\sum_{t=0}^{T-1}
\frac{g_t^{2}}{(1-\beta)\sum_{k=0}^{t}\beta^{\,t-k}g_k^{2}+\varepsilon^{2}}
\;\le\;
\frac{\alpha^{2}}{1-\beta}\,
\log\!\Bigl(1+\frac{(1-\beta)\sum_{t=0}^{T-1}g_t^{2}}{\varepsilon^{2}}\Bigr). \tag{4}
\]
\end{lemma}
\begin{proof}
Set $a_t:=g_t^{2}\ge0$ and $A_t:=\sum_{k=0}^{t}\beta^{\,t-k}a_k$.  By definition $v_{t+1}=(1-\beta)A_t$, hence
\[
\eta_t^{2}= \frac{\alpha^{2}}{(1-\beta)A_t+\varepsilon^{2}}
= \frac{\alpha^{2}}{1-\beta}\,\frac{a_t}{(1-\beta)A_t+\varepsilon^{2}}.
\]
Observe that $A_t=\beta A_{t-1}+a_t$ with $A_{-1}=0$.  For any $c>0$ the elementary inequality
\[
\frac{a_t}{\beta A_{t-1}+a_t+c}\;\le\;
\log\!\Bigl(\frac{\beta A_{t-1}+a_t+c}{\beta A_{t-1}+c}\Bigr)
\]
holds because $\log(1+x)\ge x/(1+x)$ for $x\ge0$.  Taking $c=\varepsilon^{2}/(1-\beta)$ and summing from $t=0$ to $T-1$ telescopes the logarithms:
\[
\sum_{t=0}^{T-1}\frac{a_t}{(1-\beta)A_t+\varepsilon^{2}}
\;\le\;
\frac{1}{1-\beta}\log\!\Bigl(1+\frac{(1-\beta)\sum_{t=0}^{T-1}a_t}{\varepsilon^{2}}\Bigr).
\]
Multiplying by $\alpha^{2}$ yields \eqref{eq:logsum}. \qedhere
\end{proof}

\section*{Main Proof (Step‑by‑Step Derivation)}
We prove Theorem~\ref{thm:main}.  Throughout we condition on the filtration
$\mathcal{F}_t:=\sigma(w_0,\dots,w_t,g_0,\dots,g_{t-1})$.

\subsection*{Step 1 (Smoothness descent).}
\[
\begin{aligned}
f(w_{t+1})
&\le f(w_t)+\nabla f(w_t)(w_{t+1}-w_t)+\frac{L}{2}(w_{t+1}-w_t)^2
\qquad\text{by (A1)}\\
&= f(w_t)-\eta_t\,\nabla f(w_t)g_t+\frac{L}{2}\eta_t^{2}g_t^{2}. \tag{5}
\end{aligned}
\]

\subsection*{Step 2 (Expectation handling).}
Taking total expectation of \eqref{eq:5} and using the tower property,
\[
\E[f(w_{t+1})]\le \E[f(w_t)]-\E[\eta_t\,\nabla f(w_t)g_t]
+\frac{L}{2}\E[\eta_t^{2}g_t^{2}]. \tag{6}
\]
The term $\E[\eta_t\,\nabla f(w_t)g_t]$ cannot be simplified by
replacing $g_t$ with its mean because $\eta_t$ depends on $g_t$.
We therefore write
\[
\eta_t\,\nabla f(w_t)g_t
= \eta_t\|\nabla f(w_t)\|^{2}
   +\eta_t\,\nabla f(w_t)\bigl(g_t-\nabla f(w_t)\bigr)
\]
and bound the second part using Cauchy–Schwarz and (A2):
\[
\begin{aligned}
\big|\E[\eta_t\,\nabla f(w_t)(g_t-\nabla f(w_t))]\big|
&\le \E\!\big[\eta_t\,\|\nabla f(w_t)\|\,|g_t-\nabla f(w_t)|\big]  \\
&\le \frac12\E\!\big[\eta_t\|\nabla f(w_t)\|^{2}\big]
   +\frac12\E\!\big[\eta_t (g_t-\nabla f(w_t))^{2}\big]   \\
&\le \frac12\E\!\big[\eta_t\|\nabla f(w_t)\|^{2}\big]
   +\frac12\sigma^{2}\E[\eta_t],
\end{aligned}
\]
where we used $\E[(g_t-\nabla f(w_t))^{2}\mid\mathcal{F}_t]\le\sigma^{2}$ (A2) and the fact that $\eta_t$ is $\mathcal{F}_{t+1}$‑measurable.  Consequently,
\[
\E[\eta_t\,\nabla f(w_t)g_t]
\ge \frac12\E[\eta_t\|\nabla f(w_t)\|^{2}]
      -\frac12\sigma^{2}\E[\eta_t]. \tag{7}
\]

\subsection*{Step 3 (Bounding $\E[g_t^{2}]$).}
From (A2) we have
\[
\E[g_t^{2}\mid\mathcal{F}_t]
= \|\nabla f(w_t)\|^{2}+\Var(g_t\mid\mathcal{F}_t)
\le \|\nabla f(w_t)\|^{2}+\sigma^{2}. \tag{8}
\]

\subsection*{Step 4 (Insert (7) and (8) into (6)).}
\[
\begin{aligned}
\E[f(w_{t+1})]
&\le \E[f(w_t)]
   -\frac12\E[\eta_t\|\nabla f(w_t)\|^{2}]
   +\frac12\sigma^{2}\E[\eta_t] \\
&\quad +\frac{L}{2}\E\!\big[\eta_t^{2}
      (\|\nabla f(w_t)\|^{2}+\sigma^{2})\big].
\end{aligned}
\tag{9}
\]

\subsection*{Step 5 (Upper bound on $\eta_t$).}
From Lemma~\ref{lem:vtbound} we have $v_{t+1}\le G^{2}$, hence
\[
\eta_t
= \frac{\alpha}{\sqrt{v_{t+1}}+\varepsilon}
\le \frac{\alpha}{\varepsilon}
=: \bar\eta . \tag{10}
\]
Because of the additional restriction $\alpha\le\varepsilon/(2L)$ in (A4),
\[
\frac{L}{2}\bar\eta
= \frac{L\alpha}{2\varepsilon}
\le \frac12 . \tag{11}
\]
Thus for every $t$,
\[
\frac{L}{2}\eta_t^{2}
\le \frac{L}{2}\bar\eta\,\eta_t
\le \frac12\eta_t . \tag{12}
\]
% CORRECTED: the previous proof incorrectly used $\alpha\le1/L$; the new condition $\alpha\le\varepsilon/(2L)$ guarantees (12).

\subsection*{Step 6 (Simplify (9) using (12)).}
Applying (12) to the term involving $\|\nabla f(w_t)\|^{2}$ yields
\[
\E[f(w_{t+1})]
\le \E[f(w_t)]
   -\frac12\E[\eta_t\|\nabla f(w_t)\|^{2}]
   +\frac12\sigma^{2}\E[\eta_t]
   +\frac{L\sigma^{2}}{2}\E[\eta_t^{2}]. \tag{13}
\]

\subsection*{Step 7 (Introduce the optimality gap).}
Define $\Delta_t:=\E[f(w_t)]-f^{*}\ge0$.  Taking expectation in (13) gives
\[
\Delta_{t+1}
\le \Delta_{t}
-\frac12\,\E[\eta_t\|\nabla f(w_t)\|^{2}]
+\frac12\sigma^{2}\E[\eta_t]
+\frac{L\sigma^{2}}{2}\E[\eta_t^{2}]. \tag{14}
\]

\subsection*{Step 8 (Lower bound $\eta_t$).}
From Lemma~\ref{lem:eta} and the bound $|g_t|\le G$ (A3) we obtain a deterministic
lower bound
\[
\eta_t\;\ge\;\frac{\alpha}{\sqrt{1-\beta}\,G+\varepsilon}
=: \underline\eta . \tag{15}
\]

\subsection*{Step 9 (Upper bound $\eta_t^{2}$ – logarithmic sum).}
Lemma~\ref{lem:logsum} with $a_t=g_t^{2}$ gives
\[
\sum_{t=0}^{T-1}\E[\eta_t^{2}]
\;\le\;
\frac{\alpha^{2}}{1-\beta}\,
\log\!\Bigl(1+\frac{(1-\beta)\sum_{t=0}^{T-1}\E[g_t^{2}]}{\varepsilon^{2}}\Bigr)
\;\le\;
\frac{\alpha^{2}}{1-\beta}\,
\log\!\Bigl(1+\frac{(1-\beta)G^{2}T}{\varepsilon^{2}}\Bigr). \tag{16}
\]
% ADDED: rigorous derivation of the logarithmic bound that was missing in the original proof.

\subsection*{Step 10 (Summation of the descent inequality).}
Summing \eqref{eq:14} over $t=0,\dots,T-1$ and discarding the non‑negative term $\Delta_T$ yields
\[
\frac12\underline\eta\sum_{t=0}^{T-1}\E\!\big[\|\nabla f(w_t)\|^{2}\big]
\le \Delta_{0}
+\frac12\sigma^{2}\sum_{t=0}^{T-1}\E[\eta_t]
+\frac{L\sigma^{2}}{2}\sum_{t=0}^{T-1}\E[\eta_t^{2}]. \tag{17}
\]

\subsection*{Step 11 (Bounding $\sum\eta_t$).}
Since $\eta_t\le\bar\eta=\alpha/\varepsilon$, a crude bound suffices:
\[
\sum_{t=0}^{T-1}\E[\eta_t]\le \frac{\alpha}{\varepsilon}\,T . \tag{18}
\]

\subsection*{Step 12 (Insert the bounds (15), (16) and (18) into (17)).}
Using $\underline\eta$ from \eqref{eq:15} and the estimates \eqref{eq:16}–\eqref{eq:18} we obtain
\[
\begin{aligned}
\sum_{t=0}^{T-1}\E\!\big[\|\nabla f(w_t)\|^{2}\big]
&\le
\frac{2\Delta_{0}}{\underline\eta}
   +\frac{\sigma^{2}\alpha}{\varepsilon\underline\eta}\,T
   +\frac{L\sigma^{2}\alpha^{2}}{(1-\beta)\underline\eta}
     \log\!\Bigl(1+\frac{(1-\beta)G^{2}T}{\varepsilon^{2}}\Bigr) .
\end{aligned}
\tag{19}
\]
Dividing by $T$ and recalling $\Delta_{0}=f(w_0)-f^{*}$ gives exactly the bound \eqref{eq:1} stated in Theorem~\ref{thm:main}.  The $O(\log T/\sqrt{T})$ rate follows by choosing a constant $\alpha$ (allowed by (A4)) and observing that the logarithmic term grows slower than $\sqrt{T}$.  \qed

\section*{Rate Derivation (With Constants)}
Setting $\alpha$ to any constant satisfying $\alpha\le\varepsilon/(2L)$,
\[
\frac{1}{T}\sum_{t=0}^{T-1}\E[\|\nabla f(w_t)\|^{2}]
\;\le\;
\underbrace{\frac{2\big(f(w_0)-f^{*}\big)}{\alpha}}_{C_{1}}\frac{1}{T}
\;+\;
\underbrace{\frac{2L\sigma^{2}\alpha}{1-\beta}}_{C_{2}}
\frac{\log\!\big(1+O(T)\big)}{T}
\;=\;
O\!\bigg(\frac{\log T}{\sqrt{T}}\bigg).
\]

\section*{Special Cases}
\begin{itemize}
\item \textbf{No averaging ($\beta=0$).}  The bound reduces to the classic SGD rate $O(1/\sqrt{T})$ because the logarithmic factor becomes $\log(1+G^{2}T/\varepsilon^{2})=O(\log T)$ and the prefactor $1/(1-\beta)$ equals $1$.
\item \textbf{Deterministic gradients ($\sigma=0$).}  The variance term disappears and we obtain a pure $O(1/T)$ decay of the average squared gradient norm.
\item \textbf{Large $\varepsilon$.}  If $\varepsilon\gg G$, then $\eta_t\approx\alpha/\varepsilon$ is essentially constant, and RMSProp behaves like SGD with a fixed stepsize, recovering the standard $O(1/\sqrt{T})$ bound.
\end{itemize}

\begin{thebibliography}{9}
\bibitem{rmsprop_analysis}
T. Tieleman and G. Hinton, ``Lecture 6.5—RMSProp: Divide the gradient by a running average of its recent magnitude,'' \emph{COURSERA: Neural Networks for Machine Learning}, 2012.
\end{thebibliography}

\end{document}

% ===== CORRECTION SUMMARY =====
% 1. [Step 2] Issue: η_t depends on g_t, so unbiasedness cannot be applied directly.
%    Fix: Decomposed the product, used Cauchy–Schwarz and variance bound to obtain inequality (7).
% 2. [Step 5] Issue: Inequality (9) required ε≥1, which was not assumed.
%    Fix: Added the extra condition α ≤ ε/(2L) (A4) and derived a valid bound (12).
% 3. [Step 6] Issue: Relied on the invalid inequality from Step 5.
%    Fix: Re‑derived the descent inequality (13) using the correct bound (12).
% 4. [Step 13] Issue: Missing proof of the logarithmic bound (18).
%    Fix: Introduced Lemma 3 with a full telescoping‑log proof and used it in (16)–(19).
% 5. Updated theorem statement to reflect the new hyper‑parameter requirement.